This section introduces several lemmas, definitions, and conventions that we use in the main results.
We use the notation for multifold sums introduced by Donald Knuth in~\cite{knuth1993johann}.
\begin{definition}[Multifold Sums of Powers]
  \label{def:multifold-sums-of-powers-recurrence}
  For non-negative integers $r$, $n$, and $m$, we define \emph{multifold sums of powers} by
  \begin{align*}
    \KnuthRFoldSum{0}{n}{m}   &= n^m, \\
    \KnuthRFoldSum{1}{n}{m}   &= \KnuthRFoldSum{0}{1}{m}
    + \KnuthRFoldSum{0}{2}{m} + \cdots + \KnuthRFoldSum{0}{n}{m}, \\
    \KnuthRFoldSum{r+1}{n}{m} &= \KnuthRFoldSum{r}{1}{m}
    + \KnuthRFoldSum{r}{2}{m} + \cdots + \KnuthRFoldSum{r}{n}{m}.
  \end{align*}
\end{definition}

We use the following notation for the \emph{falling factorial}.
\begin{definition}[Falling Factorials]
  \label{def:falling-factorials}
  For integers $n$ and $k$,
  \begin{align*}
    \fallingFactorial{n}{k} =
    \begin{cases}
      0, \quad & \mathrm{if \; } k < 0, \\
      1, \quad & \mathrm{if \; } k = 0, \\
      \prod_{j=0}^{k-1} (n-j), \quad & \mathrm{if \; } k>0.
    \end{cases}
  \end{align*}
\end{definition}
J. F. Steffensen discusses generalized factorials in~\cite{steffensen1933definition}.

\begin{lemma}[Forward Newton's formula]
  \label{lem:forward-newtons-formula}
  Let $f\colon \mathbb{Z} \to \mathbb{C}$ be a function and let $t \in \mathbb{Z}$. For every $n \in \mathbb{Z}$,
  \begin{align*}
    f(n)
    = \sum_{k=0}^{\infty} \frac{\fallingFactorial{n-t}{k}}{k!} \Delta^{k} f(t)
    = \sum_{k=0}^{\infty} \binom{n-t}{k} \Delta^{k} f(t),
  \end{align*}
  where the $k$-th \emph{forward difference} $\Delta^{k} f(t)$ of $f$
  evaluated at $t$ is defined by
  \begin{align*}
    \Delta^{k} f(t)
    = \tsum_{j=0}^{k} (-1)^{k-j} \tbinom{k}{j} f(t+j).
  \end{align*}
  \begin{proof}
    Newton introduced this formula in~\cite{newton1850newton}. Steffensen expressed it in modern notation in~\cite[Section 2, Equation (19)]{steffensen1927interpolation}. The identity
    $\frac{\fallingFactorial{n-t}{k}}{k!} = \tbinom{n-t}{k}$
    gives the second expression.
  \end{proof}
\end{lemma}

The power case follows directly from the preceding formula.
\begin{lemma}[Forward Newton's formula for powers]
  \label{lem:forward-newtons-formula-for-powers}
  For non-negative integers $n$ and $m$ and an arbitrary integer $t$,
  \begin{align*}
    n^m = \sum_{k=0}^{m} \binom{n-t}{k} \Delta^{k} t^m.
  \end{align*}
  \begin{proof}
    This is the special case of Lemma~\eqref{lem:forward-newtons-formula} obtained by taking $f(n)=n^m$.
  \end{proof}
\end{lemma}

We next recall a generalized hockey-stick identity.
\begin{lemma}[Generalized hockey-stick identity]
  \label{lem:generalized-hockey-stick-identity}
  For arbitrary integers $a$, $b$, and $j$,
  \begin{align*}
    \sum_{k=a}^{b} \binom{k}{j} = \binom{b+1}{j+1} - \binom{a}{j+1}.
  \end{align*}
  \begin{proof}
    We have
    $\tsum_{k=a}^{b} \tbinom{k}{j}
    = \tbinom{a}{j} + \tbinom{a+1}{j} + \cdots + \tbinom{b}{j}
    = \left( \tsum_{k=0}^{b} \tbinom{k}{j} \right) - \left( \tsum_{k=0}^{a-1} \tbinom{k}{j} \right)$.
    The hockey-stick identity $\tsum_{k=0}^{n} \tbinom{k}{j} = \tbinom{n+1}{j+1}$ therefore gives
    \begin{align*}
      \tsum_{k=a}^{b} \tbinom{k}{j}
      = \left( \tsum_{k=0}^{b} \tbinom{k}{j} \right) - \left( \tsum_{k=0}^{a-1} \tbinom{k}{j} \right)
      = \tbinom{b+1}{j+1} - \tbinom{a}{j+1}.
    \end{align*}
    This proves the identity.
  \end{proof}
\end{lemma}

The following specialization will be useful in deriving the main results.
\begin{lemma}[Forward hockey-stick identity]
  \label{lem:forward-hockey-stick-identity}
  For an integer $n\geq1$ and arbitrary integers $t$, $k$, and $r$,
  \begin{align*}
    \sum_{j=1}^{n} \binom{j-t+r}{k+r}
    = \binom{n-t+r+1}{k+r+1} - \binom{1-t+r}{k+r+1}.
  \end{align*}
  \begin{proof}
    Set $a=1-t+r$ and $b=n-t+1$ in Lemma~\eqref{lem:generalized-hockey-stick-identity}.
  \end{proof}
\end{lemma}

We also record a formula for multifold sums of zero powers.
\begin{lemma}[Multifold sums of zero powers]
  \label{lem:multifold-sum-of-zero-powers}
  For non-negative integers $r$ and $n$,
  \begin{align*}
    \KnuthRFoldSum{r}{n}{0} = \binom{r+n-1}{r}.
  \end{align*}
  \begin{proof}
    \begin{enumerate}
      \item When $r=0$, the definition gives $\KnuthRFoldSum{0}{n}{0} = 1$.
      \item When $r=1$,
        $\KnuthRFoldSum{1}{n}{0}
        = \tsum_{k=1}^{n} 1 = \tbinom{n}{1}$.
      \item When $r=2$,
        $\KnuthRFoldSum{2}{n}{0}
        = \tsum_{k=1}^{n} \tbinom{k}{1}
        = \tbinom{n+1}{2}$.
      \item When $r=3$,
        $\KnuthRFoldSum{3}{n}{0}
        = \tsum_{k=1}^{n} \tbinom{k+1}{2}
        = \tbinom{n+2}{3}$.
      \item For general $r$, the recurrence and hockey-stick identity give
        \begin{align*}
          \KnuthRFoldSum{r}{n}{0}
          &= \tsum_{k=1}^{n} \KnuthRFoldSum{r-1}{k}{0}
           = \tsum_{k=1}^{n} \tbinom{k+r-2}{r-1} \\
          &= \tbinom{r+n-1}{r}.
        \end{align*}
    \end{enumerate}
    This proves the formula.
  \end{proof}
\end{lemma}

The next lemma gives the upper-index negation identity for binomial coefficients.
\begin{lemma}[Upper Binomial negation]
  \label{lem:upper-binomial-negation}
  For integers $n$ and $k$,
  \begin{align*}
    \binom{-n}{k} = (-1)^{k} \binom{k+n-1}{k}.
  \end{align*}
  \begin{proof}
    This identity is given in~\cite[p. 164, Equation~(5.14)]{graham1994concrete}.
  \end{proof}
\end{lemma}

Finally, we adopt the \emph{zero-power convention}.
\begin{convention}[Zero Powers]
  For every $x$,
  \begin{align*}
    x^0 = 1.
  \end{align*}
\end{convention}
Donald Knuth discusses the convention $x^0=1$, including the case $x=0$, in \emph{Concrete Mathematics}~\cite[p. 162]{graham1994concrete} and in~\cite{knuth1992notesnotation}.
